% Part: computability
% Chapter: computability-theory
% Section: prop-reduce

\documentclass[../../../include/open-logic-section]{subfiles}

\begin{document}

\olfileid[pt]{cmp}{thy}{ppr}
\olsection{Propriedades da redutibilidade}

Escrevemos $A \leq_m B$ se $A$ se reduz a~$B$, e essa notação sugere
que, se $A \leq_m B$, então $A$ não é ``mais difícil que''~$B$, e
$B$ é ``tão difícil quanto ou mais difícil que''~$A$, e que
$\le_m$, assim como a relação usual $\le$ sobre os números, ordena os
conjuntos (ou problemas de decisão) por sua complexidade. As duas
proposições a seguir sustentam essa intuição. A primeira diz que
$\le_m$ é transitiva.

\begin{prop}
\ollabel{prop:trans-red}
Se $A \leq_m B$ e $B \leq_m C$, então $A \leq_m C$.
\end{prop}

\begin{proof}
Compor uma redução de $A$ a~$B$ com uma redução de $B$ a~$C$ produz
uma redução de $A$ a~$C$.
\end{proof}

\begin{prob}
Demonstre \olref{prop:trans-red} mostrando que, se $f$ e~$g$ são,
respectivamente, reduções muitos-para-um de $A$ a~$B$ e de $B$ a~$C$,
então $\comp{f}{g}$ é uma redução muitos-para-um de $A$ a~$C$.
\end{prob}

\begin{prop}
\ollabel{prop:reduce}
Sejam $A$ e $B$ conjuntos quaisquer, e suponha que $A \leq_m B$.
\begin{enumerate}
\item Se $B$ é computavelmente enumerável, $A$ também é.
\item Se $B$ é computável, $A$ também é.
\end{enumerate}
\end{prop}

\begin{proof}
Seja $f$ uma redução muitos-para-um de $A$ a~$B$. Para a primeira
afirmação, basta verificar que, se $B$ é o domínio de uma função
parcial~$g$, então $A$ é o domínio de~$\comp{f}{g}$:
\begin{align*}
x \in A & \text{ se, e somente se, } f(x) \in B \\
& \text{ se, e somente se, }  g(f(x)) \fdefined.
\end{align*}

Para a segunda afirmação, lembre-se de que, se $B$~é computável, então
$B$ e~$\Complement{B}$ são computavelmente enumeráveis
(\olref[cmp]{thm:ce-comp}). Não é difícil verificar que $f$~também é
uma redução muitos-para-um de $\Complement{A}$ a $\Complement{B}$;
portanto, pela primeira parte desta demonstração, $A$
e~$\Complement{A}$ são !!{computably
enumerable}. Assim, $A$ também é
computável por \olref[cmp]{thm:ce-comp}. (Como alternativa, você pode
verificar que $\Char{A} = \comp{f}{\Char{B}}$; portanto, se
$\Char{B}$ é computável, $\Char{A}$ também é.)
\end{proof}

\begin{prob}
Suponha que $f$ seja uma redução muitos-para-um de $A$ a~$B$. Mostre
que $f$ também é uma redução muitos-para-um de $\Complement{A}$ a
$\Complement{B}$.
\end{prob}

\begin{prob}
Mostre que, se $f\colon A \to B$ é uma redução muitos-para-um, então
$\Char{A} = \comp{f}{\Char{B}}$.
\end{prob}

\begin{digress}
Uma noção mais geral de redutibilidade, chamada \emph{redutibilidade
de Turing}, é útil em outros contextos, sobretudo para demonstrar
resultados de indecidibilidade. Observe que, por
\olref[cmp]{cor:comp-k}, o complemento de~$K_0$ não é redutível a~$K_0$,
pois não é computavelmente enumerável. Intuitivamente, porém, se você
soubesse as respostas a perguntas sobre $K_0$, também saberia as
respostas a perguntas sobre seu complemento. Diz-se que um conjunto
$A$ é redutível por Turing a~$B$ se for possível determinar respostas
a perguntas sobre~$A$ usando um procedimento computável que possa
fazer perguntas sobre~$B$. Essa noção é mais liberal do que a
redutibilidade muitos-para-um, na qual (1)~só se permite fazer uma
pergunta sobre $B$, e (2)~uma resposta ``sim'' precisa se traduzir em uma
resposta ``sim'' à pergunta sobre $A$, e de modo semelhante para
``não''. Ainda é verdade que, se $A$~é redutível por Turing a~$B$ e
$B$~é computável, então $A$~também é computável (embora, como vimos,
o enunciado análogo não valha para a enumerabilidade computável).

Você deve refletir sobre as várias noções de redutibilidade que
discutimos e compreender as distinções entre elas. Neste capítulo,
porém, trataremos apenas da redutibilidade muitos-para-um. A propósito,
os dois tipos de redutibilidade discutidos no parágrafo anterior têm
análogos na complexidade computacional, com a exigência adicional de
que as máquinas de Turing sejam executadas em tempo polinomial: a
versão da redutibilidade muitos-para-um para a complexidade é conhecida
como \emph{redutibilidade de Karp}, enquanto a versão da redutibilidade
de Turing para a complexidade é conhecida como \emph{redutibilidade de
Cook}.
\end{digress}

\end{document}
